\section{Conclusion}
\label{sec:conclusion}

This paper presents a robust harmonic hosting capacity method that combines the uncertain angles of the harmonic currents, treated in probability through chance constraints, with uncertain cable and overhead-line parameters, treated in the worst case over the intervals of their physical inputs. The parameters follow at every harmonic frequency from the physical construction of the cables and lines, and the zero sequence from the windings and earthing of the transformers. With a constant risk factor on exact moments, the robust counterpart is a convex semi-infinite second-order cone program, which a cutting-set method solves with upper and lower bounds. Sign-predicted corners make its worst-case analysis reliable in high dimension, where random sampling misses the worst cases.

On the IEEE 30-bus system with underground cables and overhead lines, the uncertainty of the network costs 14--23\,\% of the hosting capacity up to the 25th harmonic. The cost is uneven over the harmonics, with robust-to-nominal ratios from 0.45 to 0.99, so a robust budget is a new allocation rather than a uniform derating of the nominal one. The worst cases lie at the least damped network, i.e., the coldest cables and lines and the lowest earth resistivity, near the parallel resonances, which the line charging moves and multiplies. Up to the 50th harmonic without recombination, the earthing of the 33\,kV transformers changes the zero-sequence hosting capacity by up to 42\,\% and leaves the positive and negative sequence untouched, and its effect fades at the higher orders, where the zero-sequence current returns through the capacitances of the cables and lines. The Gaussian factor on exact moments stays within 1\,\% of the exact quantile, and its violation probability at the worst cases is at most 1.3 times the risk level. Every case solves within minutes.

These results question the common practice of extrapolating the fundamental-frequency parameters to the harmonics. Scaling the resistance with~$\sqrt{\Harmonic}$ and the reactance with~$\Harmonic$ overestimates the damping of the lines by a factor of two to three and the reactance of the cables by up to 2.4, misplaces the resonance next to the binding limit, and yields an allocation that violates the limits of the physical model, even when made robust over factor intervals derived from the same physical inputs. In the zero sequence, the rule fails altogether, as the resistance follows the return path through the cable screens or the earth rather than the skin effect of the conductor. The harmonic impedance of cables and lines therefore has to be computed from their physical inputs, and the uncertainty of these inputs belongs in the hosting capacity problem itself.